\section{Fording at the level of formulas: design questions settled}
\label{sec:formulas}

The model $\ML$ was built once; this section spends it.  Each subsection
takes a design question documented in implementations of related
translations --- where may the forded inversion disjunction be inlined?
which guards are load-bearing?  do match equations need index premises?
what may be said about indices without saying it about type formers? ---
and settles it as a theorem about $\ML$ or about the theory $\ThE$.
Throughout, $\Env$ is well formed and definitional with $\Sfrag$
declarations.

\subsection{Per-occurrence expansion of forded guards}
\label{sec:formulas-expansion}

The translation of \cref{sec:translation} keeps guard atoms atomic and
states the ford once per family, in the inversion axiom.  A shallower
variant --- the analogue, for families, of the inline expansion of
refinement payloads --- replaces a guard atom, at any chosen occurrence,
by the corresponding disjunction of constructor payloads and index
equations.  This is sound, in both polarities, and changes nothing
logically:

\begin{definition}[One-level expansion]
\label{def:expansion}
For a datatype $I$ as in \cref{def:theory}, define the first-order
formula (with free variables $w, \tup{z}$)
\[
  \Xp{I}(w, \tup{z}) \;\defeq\;
  \bigvee_i \exists \tup{y}.\
  \textstyle\bigwedge_j \guard{A_{ij}}{y_j} \wedge
  \bigwedge_l \F{\varphi_{il}} \wedge
  w \feq \Ihat{\wcns{c}}_i(\tup{y}) \wedge
  \textstyle\bigwedge_{h} z_h \feq \compf(\erase{t_{ih}}),
\]
i.e.\ the right-hand side of the inversion axiom; for an inductive
predicate $J$, let $\Xp{J}(\tup{z})$ be the right-hand side of its
case-analysis axiom.  An \emph{expansion variant} of the translation is
any pair $(\guardf', \Ff')$ obtained from \cref{def:guard-F} by
replacing, at any chosen set of occurrences in the structural recursion,
the clause value $\Ihat{I}(u, \compf(\erase{\tup{a}}))$ by
$\Xp{I}(u, \compf(\erase{\tup{a}}))$ and/or atoms
$\Ihat{J}(\compf(\erase{\tup{a}}))$ by
$\Xp{J}(\compf(\erase{\tup{a}}))$ (arguments substituted for $w,
\tup{z}$).  The variant theory $\Th'(\Env, G)$ is \cref{def:theory} with
items 3--5 (specification, premise and goal formulas and their auxiliary
axioms) computed by $(\guardf', \Ff')$; items 1--2 (structural axioms)
are unchanged.
\end{definition}

Note that the expansion is \emph{one level}: the guards
$\guard{A_{ij}}{y_j}$ inside $\Xp{I}$ are the atomic guards of the base
translation, so the definition is not recursive and expanded formulas
are finite regardless of recursion in $I$.  Constraint discharge remains
pay-as-you-go, one constructor layer per expansion, exactly as in the
source calculus (\cref{sec:fording-fording}).

\begin{lemma}[Expansion is valid]
\label{lem:expansion-valid}
$\ML \vDash \forall w\,\tup{z}.\ \Ihat{I}(w, \tup{z}) \leftrightarrow
\Xp{I}(w, \tup{z})$, and $\ML \vDash \forall \tup{z}.\
\Ihat{J}(\tup{z}) \leftrightarrow \Xp{J}(\tup{z})$.
\end{lemma}

\begin{proof}
Left to right is the validity of inversion resp.\ case analysis
(\cref{thm:env-validity}).  Right to left, for $I$: given a disjunct's
witnesses $\tup{d}$ satisfying its conjuncts, \cref{lem:coincidence}
converts the guard and premise conjuncts into memberships and truths at
$\rho_0 = [\tup{y} \mapsto \tup{d}\,]$;
\cref{lem:index-reflection}(5) gives $(\sem{\tup{t}_i}_{\rho_0},
\Ihat{\wcns{c}}_i^{\ML}(\tup{d})) \in \semr{I}{}$; and the equation
conjuncts identify $\den{w}{}$ with the constructor value and
$\den{\tup{z}}{}$ with the index tuple (using \cref{lem:comp-adequacy}
for the right-hand sides).  For $J$: the stage construction closes
$\pv{J}$ under exactly the displayed conditions.
\end{proof}

\begin{theorem}[Soundness of expansion variants]
\label{thm:expansion}
Let $(\guardf', \Ff')$ be any expansion variant.  Then
\cref{lem:coincidence} holds verbatim for $(\guardf', \Ff')$;
consequently $\ML \vDash \Th'(\Env,G)$, and if $\Th'(\Env,G)
\provesfol \Ff'(G)$ then $\pv{G} = 1$.  All corollaries of
\cref{sec:soundness-main} hold for the variant.
\end{theorem}

\begin{proof}
The coincidence proof is by structural induction on types and
propositions; the only change is at a replaced occurrence, where
$\ML, \rho \vDash \Xp{I}(u, \compf(\erase{\tup{a}}))$ iff $\ML, \rho
\vDash \Ihat{I}(u, \compf(\erase{\tup{a}}))$ by
\cref{lem:expansion-valid} (a universally valid equivalence may be
instantiated at the argument terms), and the base-translation
coincidence handles the atom.  The auxiliary axioms emitted by
computing the $\compf$-translations inside expanded occurrences are
compilation axioms, valid by \cref{lem:compile-sound}.  With
coincidence in place, the proofs of \cref{lem:spec-valid},
\cref{thm:env-validity} and \cref{thm:soundness} go through unchanged
for items 3--5, and items 1--2 are shared with $\ThE$.
\end{proof}

\begin{proposition}[Expansion changes nothing logically]
\label{prop:expansion-derivable}
$\ThE \provesfol \forall w \tup{z}.\ \Ihat{I}(w,\tup{z})
\leftrightarrow \Xp{I}(w,\tup{z})$ (from the constructor-typing and
inversion axioms of $I$), and likewise for $J$ from clause and
case-analysis axioms.  Consequently $\ThE$ and any $\Th'(\Env,G)$ are
mutually derivable, and the choice among expansion variants is purely
operational.
\end{proposition}

\begin{proof}
Left to right is the inversion axiom.  Right to left: in the $i$-th
disjunct, the witnesses' guards and premises feed the constructor-typing
axiom of $\wcns{c}_i$, yielding
$\Ihat{I}(\Ihat{\wcns{c}}_i(\tup{y}), \compf(\erase{\tup{t}_i}))$;
rewriting with $w \feq \Ihat{\wcns{c}}_i(\tup{y})$ and $z_h \feq
\compf(\erase{t_{ih}})$ gives $\Ihat{I}(w, \tup{z})$.  Mutual
derivability of the theories: each axiom of $\Th'$ differs from the
corresponding axiom of $\ThE$ by replacements of subformulas by
$\ThE$-provably equivalent ones (both theories contain items 1--2), so
each direction follows by the replacement theorem of classical
first-order logic.
\end{proof}

\begin{example}[The $\mathsf{reflect}$ payoff]
\label{ex:reflect-expansion}
For $\mathsf{rfl}_{\varphi_0}$ (\cref{ex:running-axioms}), a hypothesis
occurrence $r : \mathsf{rfl}_{\varphi_0}(b)$ inside a premise or goal
expands to
\[
  \bigl(\F{\varphi_0} \wedge r \feq \Ihat{\wcns{RT}} \wedge b \feq
  \Ihat{\cstrue}\bigr) \vee
  \bigl(\F{\neg\varphi_0} \wedge r \feq \Ihat{\wcns{RF}} \wedge b \feq
  \Ihat{\csfalse}\bigr).
\]
Both directions of the view are now local: from $b \feq \Ihat{\cstrue}$
the second disjunct dies by discrimination on $\bool$, leaving
$\F{\varphi_0}$ (elimination is no longer too weak); and in a positive
position, establishing the instance at tag $\Ihat{\wcns{RF}}$ together
with $b \feq \Ihat{\cstrue}$ is impossible for the same reason
(introduction cannot cheat the index).  In the atomic variant the same
two inferences route through the declaration-level inversion axiom ---
one resolution step more, contingent on premise selection.  Which
variant is operationally better is an empirical question; the content
of \cref{thm:expansion,prop:expansion-derivable} is that it is
\emph{only} an operational question.
\end{example}

\begin{remark}[Existential residue]
Expansion introduces $\exists \tup{y}$ in guard positions; in negative
positions these become universals after clausification.  For
enumeration-like and subset-like families the residue is empty or a
single carrier variable after \cref{lem:n-one} ($\mathsf{rfl}$:
none; $\mathsf{bnd}$: one), which is exactly the class of families for
which implementations enable expansion; for recursive families like
$\vect$ the atomic guard is the reasonable default.  Nothing in the
soundness statement forces either choice.
\end{remark}

\subsection{Which guards are load-bearing}
\label{sec:formulas-guards}

The companion paper proved that for non-indexed datatypes the guard of
the inversion axiom is semantically mandatory (false at junk) while the
per-constructor equations need no guards at all.  For families both
halves sharpen.

\begin{proposition}[Unguarded inversion is inconsistent]
\label{prop:unguarded-inconsistent}
Let $I$ be a datatype some constructor $\wcns{c}_i$ of which has a
constructor-headed index term: $t_{ih} = \wcns{c}'(\dots)$ for some
position $h$, where $\wcns{c}'$ belongs to a datatype $I'$ that has a
second constructor $\wcns{d} \neq \wcns{c}'$.  Let $\mathrm{Inv}^-_I$
be the inversion axiom of $I$ with its guard atom $\Ihat{I}(w,
\tup{z})$ deleted.  Then $\{\mathrm{Inv}^-_I\}$ together with the
discrimination axioms of $I$ and of $I'$ --- axioms the translation
must emit in any case --- is inconsistent.  In particular, for $\vect$
(\cref{ex:running-axioms}) the theory $\{\mathrm{Inv}^-_{\vect}\} \cup
\{\text{discrimination for } \nat, \vect\}$ is inconsistent.
\end{proposition}

\begin{proof}
Instantiate the universally quantified variables of $\mathrm{Inv}^-_I$
by $w \defeq \Ihat{\wcns{c}}_i(\tup{x})$, $z_h \defeq
\Ihat{\wcns{d}}(\tup{x}')$, the remaining $z$'s by fresh variables
($\tup{x}, \tup{x}'$ fresh).  In the resulting disjunction, every
disjunct $i' \neq i$ contains the equation
$\Ihat{\wcns{c}}_i(\tup{x}) \feq \Ihat{\wcns{c}}_{i'}(\tup{y})$,
refuted by discrimination for $I$; the disjunct $i$ contains
$\Ihat{\wcns{d}}(\tup{x}') \feq \compf(\erase{t_{ih}})[\tup{y}]$, whose
right-hand side is $\Ihat{\wcns{c}'}$-headed
($\compf(\erase{\wcns{c}'(\dots)})$ is an application of
$\Ihat{\wcns{c}'}$ by \cref{def:compC}), refuted by discrimination for
$I'$.  All disjuncts refuted, the theory proves $\bot$.
For $\vect$, concretely: $w \defeq \Ihat{\csvnil}$, $z \defeq
\Ihat{\csS}(x')$ --- the $\csvnil$ disjunct dies on $\Ihat{\csS}(x')
\feq \Ihat{\csO}$, the $\csvcons$ disjunct on $\Ihat{\csvnil} \feq
\Ihat{\csvcons}(n,a,v)$.
\end{proof}

\begin{remark}
\label{rem:unguarded-general}
The hypothesis delimits the blast radius precisely.  Rigid index
patterns (the proposition) cover $\vect$, $\Fin$ and every family whose
constructors specialise an index.  Enumeration-shaped families with
bare-variable or absent index constraints are caught by the
\emph{cardinality} mechanism instead: for $\mathsf{rfl}_{\varphi_0}$,
the unguarded inversion entails $\forall w.\ w \feq \Ihat{\wcns{RT}}
\vee w \feq \Ihat{\wcns{RF}}$ (each disjunct contains the corresponding
tag equation as a conjunct), a two-element bound on the whole untyped
universe; in any environment that also declares $\nat$ --- or any
datatype with three provably distinct closed terms --- instantiating
$w$ at $\Ihat{\csO}$, $\Ihat{\csS}(\Ihat{\csO})$ and
$\Ihat{\csS}(\Ihat{\csS}(\Ihat{\csO}))$ --- pairwise distinct by
discrimination and injectivity of $\nat$ --- forces two of them equal
by pigeonhole, hence $\bot$.  This is Meng and Paulson's two-element
universe \cite[\S 2.3]{MengPaulson2008}; exhaustiveness sentences are
exactly the sentences whose guards can never be dropped
\cite{BlanchetteEtAl2016}.  For a fully generic family (every index a
bare variable, constructors carrying arguments) the unguarded axiom is
false in $\ML$ --- any junk instantiation of $w$ refutes it there ---
but need not be inconsistent; soundness fails all the same, via
\cref{lem:coincidence}.  Note the strengthening against the non-indexed
case of \cite{Extraction2026}: there, guard omission made the
disjunctive axiom false in $\ML$, and inconsistency required an
auxiliary single-constructor datatype; here the family's own index
equations or tag equations detonate it.  For the implementation this
means: an indexed inversion axiom accidentally emitted unguarded is
typically not a weak axiom but a refutation-generator, poisoning every
problem it enters (cf.\ \cref{cor:consistency}).
\end{remark}

\begin{proposition}[Index arguments of guards are mandatory]
\label{prop:index-dropping}
Consider the \emph{index-forgetting} variant of the translation: guard
predicates $\Ihat{I}^-$ of arity $1$, guard clause
$\guardf^-_{I(\tup{a})}(u) = \Ihat{I}^-(u)$, and the axioms of
\cref{def:theory} recomputed accordingly (index equations deleted from
inversion, index arguments deleted from constructor-typing
conclusions).  This variant is unsound: for $\Env$ declaring $\nat$ and
$\vect$ (over $A = \nat$) there is a closed shallow goal $G$ with
$\pv{G} = 0$ whose translation is derivable,
\begin{gather*}
  G \;\defeq\; \exists v : \vect(\csO).\ \neg\,(v =_{\vect(\csO)}
  \csvnil),\\
  \Th^-(\Env, G) \provesfol \exists v.\ \Ihat{\vect}^-(v) \wedge
  \neg\,(v \feq \Ihat{\csvnil}).
\end{gather*}
\end{proposition}

\begin{proof}
$\pv{G} = 0$: by \cref{lem:index-reflection}(1) and
\cref{cor:discrimination}, every $(\sem{\csO}, d) \in \semr{\vect}{}$
decomposes by $\csvnil$ --- the $\csvcons$ index tuple is a class
$\cls{\csS(e)} \neq \cls{\csO}$ --- so $\semr{\vect(\csO)}{} =
\{\cls{\csvnil}\}$ and the matrix is false at the only element.

Derivability: the index-forgetting constructor-typing axioms are
$\Ihat{\vect}^-(\Ihat{\csvnil})$ and $\forall n\,a\,v.\ \Ihat{\nat}(n)
\to \Ihat{\nat}(a) \to \Ihat{\vect}^-(v) \to
\Ihat{\vect}^-(\Ihat{\csvcons}(n,a,v))$.  From $\Ihat{\nat}(\Ihat{\csO})$
(constructor typing for $\nat$) instantiate $n = a \defeq \Ihat{\csO}$,
$v \defeq \Ihat{\csvnil}$ to get
$\Ihat{\vect}^-(\Ihat{\csvcons}(\Ihat{\csO}, \Ihat{\csO},
\Ihat{\csvnil}))$; discrimination for $\vect$ gives
$\neg(\Ihat{\csvcons}(\dots) \feq \Ihat{\csvnil})$; the displayed
existential follows with witness
$\Ihat{\csvcons}(\Ihat{\csO},\Ihat{\csO},\Ihat{\csvnil})$.
\end{proof}

\begin{remark}[The \fstar{} \#1542 shape]
\label{rem:1542}
\Cref{prop:index-dropping} is the indexed-family form of a documented
unsoundness.  In \fstar{} \#1542 \cite{FStar1542}, the type-indexed
extensional equality $\mathsf{feq}\ \#a$ --- ``equal pointwise at domain
$a$'' --- was encoded through the SMT solver's untyped $=$; an equality
whose \emph{meaning depends on an index} was represented by an atom that
forgot the index, and facts established at index $\mathsf{nat}$ were
consumed at index $\mathsf{int}$, proving $\mathsf{False}$.  Here,
family membership --- whose meaning depends on the index --- is
represented in the variant by an atom $\Ihat{\vect}^-(v)$ that forgets
it, and membership established at index $\csS(\csO)$ is consumed at
index $\csO$.  The polarity analysis explains where the damage occurs:
forgetting indices in \emph{negative} (premise) positions merely
strengthens hypotheses and weakens the translation; in \emph{positive}
(conclusion, existential-witness) positions it lets the prover
manufacture members of the wrong instance, as in the derivation above.
Since one guard clause serves both polarities, the indices must stay.
The translation of \cref{sec:translation} is immune for a reason worth
stating precisely: it is not that its untyped $\feq$ is somehow safer
than \fstar{}'s --- it is that \emph{no index-dependent notion of the
source is represented by an index-blind notion of the target}.  Source
equality $a =_\sigma b$ is index-blind ($\sigma$ is not consulted)
because its \emph{semantics} in $\ML$ is index-blind class equality
(\cref{def:interp}), and family membership is index-aware because its
semantics is (\cref{lem:coincidence}); each atom is exactly as
discriminating as the semantic notion it names.  The residual caveat is
the intensionality of that semantics: see \cref{rem:funext-leak}.
\end{remark}

\subsection{Match equations for indexed scrutinees}
\label{sec:formulas-equations}

\begin{proposition}[Unguarded equations, impossible branches included]
\label{prop:match-equations}
Every equation emitted by the compilation of a definition matching on a
family --- with no guard on the matched term, no guards on the
constructor-argument variables, and no index premises, and
\emph{including} the equations of branches whose index patterns clash
with the scrutinee occurrence's indices --- holds in $\ML$ under every
assignment.  In particular (\cref{ex:vhead-transl}):
\begin{gather*}
  \ML \vDash \forall n.\ \Ihat{\mathsf{vhead}}(n, \Ihat{\csvnil}) \feq
  \prfc,\\
  \ML \vDash \forall n\,k\,a\,w.\
  \Ihat{\mathsf{vhead}}(n, \Ihat{\csvcons}(k,a,w)) \feq a .
\end{gather*}
\end{proposition}

\begin{proof}
An instance of \cref{lem:compile-sound}(2).  The point of displaying it:
the $\iota$-rule of $\lbox$ fires on any constructor form with no index
side condition, so validity of the equations is a fact about the
\emph{program}, blind to the typing discipline that made the branch
impossible; the branch body is the erasure of the $\bot$-elimination,
namely $\bx$.
\end{proof}

\begin{proposition}[Index-premised equations are strictly weaker]
\label{prop:premised-weaker}
Fix the definition $\mathsf{vhead}$ and let the \emph{premised variant}
of its $\csvnil$-equation be
\[
  \mathrm{Eq}^{\pi}_{\csvnil} \;\defeq\;
  \forall n.\ \Ihat{\csS}(n) \feq \Ihat{\csO} \to
  \Ihat{\mathsf{vhead}}(n, \Ihat{\csvnil}) \feq \prfc,
\]
and analogously $\mathrm{Eq}^{\pi}_{\csvcons} = \forall n\,k\,a\,w.\
\Ihat{\csS}(n) \feq \Ihat{\csS}(k) \to \Ihat{\mathsf{vhead}}(n,
\Ihat{\csvcons}(k,a,w)) \feq a$ (premises: occurrence indices $\feq$
constructor index pattern).  Then:
\begin{enumerate}[leftmargin=2em]
\item each premised equation is derivable from its unguarded form, hence
  valid in $\ML$;
\item the unguarded $\csvnil$-equation is \emph{not} derivable from the
  theory obtained from $\ThE$ by replacing $\mathsf{vhead}$'s equations
  with their premised forms;
\item deleting the $\csvnil$-equation altogether (\emph{rigid-clash
  pruning}) preserves \cref{thm:env-validity}, hence
  \cref{thm:soundness,cor:consistency}.
\end{enumerate}
\end{proposition}

\begin{proof}
(1) Weakening.  (3) Validity of a subset of a valid axiom set.

(2) Let $\Env_0$ declare $\nat$, $A \defeq \nat$, $\vect$ and
$\mathsf{vhead}$ only, and let $\Th^\pi$ be the resulting theory with
$\mathsf{vhead}$'s two equations replaced by their premised forms.  We
build $\ML' \vDash \Th^\pi$ falsifying the unguarded
$\csvnil$-equation.  A pointwise reinterpretation of
$\Ihat{\mathsf{vhead}}$ is not available: the pointer axiom $\forall
z_1 z_2.\ \Ihat{\mathsf{vhead}}^0 \fapp z_1 \fapp z_2 \feq
\Ihat{\mathsf{vhead}}(z_1, z_2)$ ties the symbol to the applicative
structure of $\Dom$, which is fixed.  Instead, reinterpret both symbols
through an \emph{alternative witness}: let
\[
  \kappa \defeq \lambda n\,v.\ \mathsf{case}(v;\,
    \{\csvnil \Rightarrow \csO \mid \csvcons(k,a,w) \Rightarrow a\})
\]
--- the erasure of $\mathsf{vhead}$ with the dead code of the
impossible branch replaced by $\csO$ --- and let $\ML'$ agree with
$\ML$ on all symbols except $(\Ihat{\mathsf{vhead}}^0)^{\ML'} \defeq
\cls{\kappa}$ and $\Ihat{\mathsf{vhead}}^{\ML'}(d_1, d_2) \defeq
\cls{\kappa\,e_1\,e_2}$ (representatives $e_i \in d_i$).  Then:
the pointer axiom holds by construction ($\fapp$ is application);
the premised $\csvnil$-equation holds vacuously, its premise
$\Ihat{\csS}(n) \feq \Ihat{\csO}$ being false
(\cref{cor:discrimination}; constructor interpretations are unchanged);
the premised $\csvcons$-equation holds because
$\kappa\,e_1\,\csvcons(\tup{e}) \lreds e_a$, and this is where $\kappa$
agrees with $\erase{\mathsf{vhead}}$; the specification axiom
$\forall n.\ \Ihat{\nat}(n) \to \forall v.\ \Ihat{\vect}(v,
\Ihat{\csS}(n)) \to \Ihat{\nat}(\Ihat{\mathsf{vhead}}(n,v))$ constrains
the symbol only at $v \in \semr{\vect(\csS\,n)}{}$-instances, which by
\cref{lem:index-reflection}(1,3) are $\csvcons$-forms ($\cls{\csvnil}$'s
unique index tuple is $\cls{\csO}$), where $\ML'$ computes the same
values as $\ML$, so the axiom transfers from
\cref{thm:env-validity}; and all axioms not mentioning
$\mathsf{vhead}$'s symbols hold as in $\ML$.  Finally the unguarded
$\csvnil$-equation fails: $\Ihat{\mathsf{vhead}}^{\ML'}(d,
\cls{\csvnil}) = \cls{\kappa\,e\,\csvnil} = \cls{\csO} \neq \cls{\bx} =
\prfc^{\ML'}$, the last two being distinct $\lred$-normal forms
(\cref{thm:confluence}).

The moral generalises: the premised theory cannot distinguish the
extracted program from any \emph{repair of its dead code}, because
every axiom that could see the dead branch is either premised on a
refutable index equation or guarded by an uninhabited instance.
\end{proof}

\begin{proposition}[Pruned branches are semantically dead]
\label{prop:pruning}
Say two terms \emph{rigidly clash} if both are constructor-headed with
distinct constructors, or both are applications of the same constructor
and some pair of corresponding informative arguments rigidly clashes;
extend to index tuples by ``some position clashes''.  If, in a
definition $c \defeq \lambda \tup{x}.\ \tcase{x_k}{\dots}{\dots}$ (or
the fix form) with $x_k : I(\tup{v})$, the occurrence indices $\tup{v}$
rigidly clash with the index pattern $\tup{t}_i$ of constructor
$\wcns{c}_i$, then for every valuation $\rho$ of $\tup{x}$ with
$\rho(x_k) \in \semr{I(\tup{v})}{\rho}$, the value $\rho(x_k)$ is not of
the form $\cls{\wcns{c}_i(\tup{e})}$.  Consequently the pruned equation
constrains $\Ihat{c}$ only at argument tuples excluded by the guards of
$c$'s specification axiom.
\end{proposition}

\begin{proof}
First, rigid clash implies semantic distinctness at all valuations: by
induction on the definition, if $u$ and $t$ rigidly clash then
$\sem{u}_{\rho_1} \neq \sem{t}_{\rho_2}$ for all $\rho_1, \rho_2$ ---
distinct constructor heads give distinct classes
(\cref{cor:discrimination}(1)), and for equal heads, componentwise
equality of classes would contradict the inner clash by the induction
hypothesis (erasure preserves constructor heads of informative
arguments).  Now suppose $\rho(x_k) = \cls{\wcns{c}_i(\tup{e})} \in
\semr{I(\tup{v})}{\rho}$.  By \cref{lem:index-reflection}(1,3) the
unique index tuple of this element is $\sem{\tup{t}_i}_{\rho_0}$, so
$\sem{\tup{v}}_\rho = \sem{\tup{t}_i}_{\rho_0}$, contradicting semantic
distinctness at the clashing position.
\end{proof}

\begin{remark}[Choice of equation form]
\label{rem:equation-choice}
\Cref{prop:match-equations,prop:premised-weaker,prop:pruning} bound the
design space: the unguarded equations are the strongest sound form and
are what the factored architecture naturally produces; index premises
are sound but strictly weaker, buying nothing semantically (their one
virtue is to keep the equation silent at wrong-index instantiations,
a reconstruction-signal nicety); pruning is sound and removes only
equations that no guarded reasoning can reach, at the price that we do
\emph{not} claim conservativity --- a first-order derivation of a goal
may in principle route through junk instances of a pruned equation, so
pruning can only be justified operationally, not by the model.  All
three verdicts are instances of the companion paper's division of
labour: equations carry computation, valid everywhere; typing content
lives in guarded specification and inversion axioms.
\end{remark}

\subsection{Indices, elements, and injectivity of type formers}
\label{sec:formulas-injectivity}

\begin{proposition}[Index uniqueness is a first-order consequence]
\label{prop:index-unique}
For every datatype $I$,
\[
  \ThE \provesfol \forall w\,\tup{z}\,\tup{z}'.\
  \Ihat{I}(w, \tup{z}) \to \Ihat{I}(w, \tup{z}') \to
  \textstyle\bigwedge_h z_h \feq z'_h,
\]
and the formula holds in $\ML$.
\end{proposition}

\begin{proof}
Model validity is \cref{lem:index-reflection}(3).  Derivability: apply
inversion to both hypotheses, obtaining disjuncts $i$ (witnesses
$\tup{y}$) and $i'$ (witnesses $\tup{y}'$) with $w \feq
\Ihat{\wcns{c}}_i(\tup{y}) \feq \Ihat{\wcns{c}}_{i'}(\tup{y}')$.  If $i
\neq i'$, discrimination refutes the branch.  If $i = i'$, injectivity
gives $y_j \feq y'_j$ for all $j$; each index equation pair $z_h \feq
\compf(\erase{t_{ih}})[\tup{y}]$, $z'_h \feq
\compf(\erase{t_{ih}})[\tup{y}']$ then yields $z_h \feq z'_h$ by
congruence of $\feq$ (the two right-hand sides are the same first-order
term at componentwise-equal arguments).  For $r_i = 0$ the two
right-hand sides are identical closed terms.
\end{proof}

\begin{remark}[Injectivity of type formers is neither expressible nor
wanted]
\label{rem:hur}
Index uniqueness is the sound first-order shadow of the type-theoretic
intuition ``$I(\tup{u})$ and $I(\tup{u}')$ are disjoint unless $\tup{u}
= \tup{u}'$'': it speaks about \emph{elements}.  The stronger,
type-level statement --- injectivity of the type former, $I\,\tup{u} =
I\,\tup{u}' \to \tup{u} = \tup{u}'$ --- is not even a formula of
$\SigFOL$: types are not terms in this translation
(\cref{def:signature}), so there is no equality between type images to
invert.  This inexpressibility is a feature.  In type theory,
injectivity of inductive type constructors is \emph{anti-classical}:
Hur's construction \cite{Hur2010} refutes excluded middle in Agda with
injective type constructors, by a Cantor-style diagonal through an
injection $(\Set \to \Set) \to \Set$ --- so a translation validating
classical logic, as any ATP-facing translation must, cannot emit any
image of it.  Encodings that do represent type expressions as
first-order terms (types-as-terms with a $\mathsf{HasType}$ predicate)
must therefore be careful that no axiom makes the type-term constructor
injective-and-surjective in a way that imports the paradox; in the
present design the question cannot arise.  The same discipline governs
constructor discrimination: the axioms of \cref{def:theory} concern
fully applied constructor images only --- there are no partially
applied constructor symbols in $\SigFOL$ at all, constructors being
$r$-ary function symbols --- which is exactly the boundary drawn by
Counterexample~35 of \cite{CockxDevriese2018}: conflict between
\emph{under-applied} constructors is refutable under functional
extensionality (Agda issue \#1497) and must not be asserted by a
translation whose source may consistently assume it.
\end{remark}

\subsection{Singleton collapse and uniqueness of identity proofs}
\label{sec:formulas-uip}

The $\tscase{}{}{}{}$ and $\tcast{}{}{}$ cases of
\cref{lem:fundamental}, and the unconditional identity equations that
user-written transports compile to (\cref{ex:eq-pred}), are the exact
locations where the model's proof irrelevance is spent as
\emph{uniqueness of identity proofs}.  We make the accounting precise.

\begin{fact}[The semantic UIP surrogate]
\label{fact:uip}
Let $J$ be a forced singleton with clause $\wcns{k} : \forall
\tup{x}{:}\tup{\tau}.\ \tup{\varphi} \to J(\tup{u}^0)$, solution
$\jmath$, and let $J(\tup{u})$ be an occurrence with solution
substitution $\theta$.  For every valuation $\rho$ with
$\pvr{J(\tup{u})}{\rho} = 1$:
\[
  \sem{\theta(\tup{u}^0)}_\rho = \sem{\tup{u}}_\rho
  \qquad\text{and}\qquad
  \pvr{\varphi_l[\theta]}{\rho} = 1 \ \text{for all } l.
\]
\end{fact}

\begin{proof}
Extracted from the $\tscase{}{}{}{}$ case of \cref{lem:fundamental},
which proves exactly this.
\end{proof}

In type-theoretic terms, \cref{fact:uip} says: the mere \emph{truth} of
the singleton family at $\tup{u}$ collapses the motive instance at the
generic (constructor) indices with the motive instance at $\tup{u}$, and
recovers the premises --- without inspecting any proof.  Gilbert,
Cockx, Sozeau and Tabareau prove that this collapse, availed as a
definitional principle for an equality-like family under proof
irrelevance, \emph{is} UIP \cite[\S 2.2]{GilbertEtAl2019}; and the
without-K analysis of pattern matching isolates the same strength in the
unification steps of deletion and of injectivity at non-linear or
non-self-unifiable indices \cite{CockxEtAl2014,CockxDevriese2018}.  Our
model validates all of them at once, because $\pv{\cdot}$ is
proof-irrelevant and \cref{lem:index-reflection} holds; consequently:

\begin{itemize}[leftmargin=2em]
\item \emph{Deletion} costs nothing: a match on $\mathsf{eq}_\sigma(a,a)$
  (or a $\tcast{}{}{}$ along it) is validated by \cref{fact:uip} with
  $\theta(\tup{u}^0)$ and $\tup{u}$ literally equal.
\item \emph{Non-linear index patterns} cost nothing: nothing in
  \cref{def:forced-singleton} or in \cref{lem:index-reflection}
  restricts repeated variables or self-unifiability; the without-K
  side conditions exist to avoid UIP, which the model has.
\item \emph{Heterogeneity} costs nothing at the target: the index
  equations of \cref{def:theory} are pointwise untyped equations; the
  formation problem that motivates John Major equality
  \cite{McBride1999} does not exist in $\SigFOL$.
\end{itemize}

What the strength \emph{costs} is reconstruction, not soundness: a
first-order proof that uses a transport-collapse equation at a type
without decidable equality may need UIP (or an axiom-free inversion
argument where one exists \cite{Monin2010,CockxDevriese2018}) to replay
in CIC.  This is the same debt the companion paper documents for
$\tcast{}{}{}$-erasure, unchanged in kind; the premise-carrying variant
below eliminates it at the price of weaker axioms.

\begin{proposition}[Werner's guarded collapse as an emission variant]
\label{prop:guarded-scase}
Let $\Env_0$ declare $\bool$, the forced singleton
$\mathsf{eq}_{\bool}$ (\cref{ex:declarations}), and
\begin{multline*}
  c \defeq \lambda x\,y{:}\bool.\ \lambda
  e{:}\mathsf{eq}_{\bool}(x,y).\ \lambda w{:}\bool.\\
  \tscase{\mathsf{eq}_{\bool}}{e}{\tup{z}.\bool}{.\,w},
\end{multline*}
with emitted equation $\forall x\,y\,w.\ \Ihat{c}(x,y,w) \feq w$
(\cref{ex:eq-pred}, at the constant motive).  The \emph{premised
variant}
\[
  \mathrm{Eq}^\pi_c \;\defeq\; \forall x\,y\,w.\ y \feq x \to
  \Ihat{c}(x,y,w) \feq w
\]
(premise: the non-trivial index equation of the ford, $\tup{u} \feq
\theta(\tup{u}^0)$) is derivable from the unguarded equation, hence
$\ML$-valid; and the unguarded equation is not derivable from the
theory $\Th^\pi$ obtained from $\Th(\Env_0, \top)$ by replacing it with
$\mathrm{Eq}^\pi_c$.
\end{proposition}

\begin{proof}
Derivability is weakening.  For non-derivability, a witness-repair
modification of $\ML$ in the style of \cref{prop:premised-weaker} is
not available: the pointer axiom forces the behaviour of $\Ihat{c}$ to
be the applicative behaviour of some element of $\Dom$, and the
premised equation constrains that behaviour on \emph{every} diagonal
triple $(d, d, w)$, junk included --- where a $\lbox$ program that
inspects its arguments is stuck; rather than pursue a definability
analysis of $\Dom$, we exhibit a finite countermodel.  Let
$\mathfrak{N}$ have universe $\{0,1\}$ with: $\feq$ interpreted as
identity; $\Ihat{\cstrue} \defeq 0$, $\Ihat{\csfalse} \defeq 1$;
$\Ihat{\bool} \defeq \{0,1\}$;
$\Ihat{\mathsf{eq}_{\bool}} \defeq \{(0,0), (1,1)\}$; $\prfc \defeq 0$;
$u \fapp v \defeq u \oplus v$ (exclusive or); $\Ihat{c}^0 \defeq 0$;
and $\Ihat{c}(x,y,w) \defeq x \oplus y \oplus w$.  Every axiom of
$\Th^\pi$ holds:
constructor typing and inversion for $\bool$ (the universe is
$\{0,1\}$), discrimination ($0 \neq 1$), the clause and case-analysis
axioms of $\mathsf{eq}_{\bool}$ (its interpretation is the diagonal of
$\Ihat{\bool}$), the specification axiom of $c$ (its conclusion
$\Ihat{\bool}(\Ihat{c}(x,y,w))$ is unconditionally true), the pointer
axiom ($((\Ihat{c}^0 \oplus x) \oplus y) \oplus w = x \oplus y \oplus
w$), and the premised equation ($y = x$ gives $x \oplus x \oplus w =
w$).  The unguarded equation fails: $\Ihat{c}(0,1,0) = 1 \neq 0$.
\end{proof}

This is precisely the situation Werner analysed for proof-irrelevant
type theories \cite{Werner2008}: the unconditional singleton reduction
$(\mathsf{Eq\_rec}\ A\ P\ a\ b\ p\ e) \rhd p$ ``easily breaks the
subject reduction property in incoherent contexts,'' and his repair
keeps the index condition as a guard: reduce only \emph{if} $a
=_{\beta\varepsilon} b$ (generalized to arbitrary singletons: if
$\tup{u} =_{\beta\varepsilon} \tup{a}$).  The premised equation is that
rule with ``provable'' in place of ``definitionally checked'' --- a
strictly weaker demand, which is why both variants are sound here while
Werner needs the check for decidability and subject reduction.  An
implementation can therefore expose the choice as a dial: unguarded
collapse (strongest, UIP-debt at reconstruction) or premised collapse
(weaker, replayable by plain inversion), per definition or globally.
Brady, McBride and McKinna's run-time warning --- collapsing ``runs the
risk of reducing a proof of something which cannot be constructed, such
as $5 \le 4$'' in open contexts \cite{BradyEtAl2004} --- is the same
point seen from the operational side; in the model-theoretic setting
the absurd-context hazard is neutralised because equations are
consequences of one consistent structure (\cref{cor:consistency}), and
what remains of it is exactly the premised-versus-unguarded choice.

\begin{remark}[Intensionality: the residual caveat]
\label{rem:funext-leak}
$\ML$ interprets equality as $\lbox$-convertibility, which is
intensional: for the standard recursive $\mathsf{add}$,
$\cls{\lambda x.\,\mathsf{add}\ \csO\ x} \neq \cls{\lambda x.\,x}$ in
$\Dom$ (the left representative is stuck at the variable) although the
two classes agree applicatively on all of $\semr{\nat}{}$.  So functional extensionality is \emph{not} validated:
environments whose premise lemmas are proved from a funext postulate
step outside \cref{thm:soundness} (\cref{rem:axioms}), though
\cref{cor:consistency} is unaffected since no extensionality axiom is
ever emitted.  This caveat is the honest residue of the \fstar{}
\#1542 analysis (\cref{rem:1542}): that incident needed \emph{both} an
extensionality axiom \emph{and} an index-blind equality encoding; the
present translation emits neither, and each half of the immunity is a
theorem --- \cref{prop:index-dropping} for the index-blindness half,
the present remark for the extensionality half.
\end{remark}
